% ==============================================================
\section{Connections to Information Theory and Extreme Values}
\label{sec:ac-advanced}
% ==============================================================

This section develops speculative connections between the CES framework and three areas of mathematics: information theory (Kullback-Leibler divergence), extreme value theory (the Gnedenko-Fisher-Tippett theorem), and large deviations (rare event probabilities). These connections suggest that the CES structure may be more fundamental than its current applications in economics indicate---it emerges naturally from information-theoretic principles and from the statistics of extreme outcomes.

\begin{calloutcaution}[Speculative Material]
The results in this section range from established theorems applied in new contexts to conjectures and open questions. We aim to sketch a research agenda rather than present finished theory. The connections are suggestive and merit further investigation.
\end{calloutcaution}

% --------------------------------------------------------------
\subsection{The Kullback-Leibler Divergence}
\label{subsec:ac-adv-kl}
% --------------------------------------------------------------

\paragraph{Definition.}
The Kullback-Leibler (KL) divergence from distribution $\bq = (q_1, \ldots, q_n)$ to distribution $\bp = (p_1, \ldots, p_n)$ is:
\begin{equation}
D_{KL}(\bp \| \bq) = \sum_{i=1}^{n} p_i \log \frac{p_i}{q_i}.
\label{eq:ac-adv-kl-def}
\end{equation}

The KL divergence measures the ``information cost'' of using $\bq$ to approximate $\bp$. It is non-negative, equals zero if and only if $\bp = \bq$, and is not symmetric.

\paragraph{Connection to CES Optimization.}
Consider the following variational problem: find the distribution $\bp$ that maximizes expected utility subject to an information constraint:
\begin{equation}
\max_{\bp} \; \sum_{i=1}^{n} p_i u_i \quad \text{subject to} \quad D_{KL}(\bp \| \bq) \leq \kappa, \quad \sum_i p_i = 1.
\label{eq:ac-adv-kl-problem}
\end{equation}

Here $u_i$ is the utility in state $i$, $\bq$ is a reference (prior) distribution, and $\kappa$ is the information budget.

\paragraph{The Solution.}
The Lagrangian is:
\begin{equation}
\mathcal{L} = \sum_i p_i u_i - \lambda \left( \sum_i p_i \log \frac{p_i}{q_i} - \kappa \right) - \mu \left( \sum_i p_i - 1 \right).
\label{eq:ac-adv-kl-lagrangian}
\end{equation}

The FOC for $p_i$ is:
\begin{equation}
u_i - \lambda \left( \log \frac{p_i}{q_i} + 1 \right) - \mu = 0.
\label{eq:ac-adv-kl-foc}
\end{equation}

Solving:
\begin{equation}
p_i = q_i \exp\left( \frac{u_i - \mu - \lambda}{\lambda} \right) = q_i \cdot \frac{\exp(u_i / \lambda)}{Z},
\label{eq:ac-adv-kl-solution}
\end{equation}
where $Z = \sum_j q_j \exp(u_j / \lambda)$ is the normalizing constant.

\begin{calloutimportant}[The Exponential Tilting]
The optimal distribution is an \emph{exponential tilting} of the prior:
\[
p_i^* \propto q_i \exp(u_i / \lambda).
\]
This is the Gibbs distribution from statistical mechanics, with $\lambda$ playing the role of temperature. Higher utility states receive more probability mass, with the tilting controlled by the information cost $\lambda$.
\end{calloutimportant}

\paragraph{The Value Function.}
The maximized objective, evaluated at the optimum, is:
\begin{equation}
V^* = \lambda \log \sum_i q_i \exp(u_i / \lambda) + \lambda \kappa.
\label{eq:ac-adv-kl-value}
\end{equation}

The term $\lambda \log \sum_i q_i \exp(u_i / \lambda)$ is the \emph{log-sum-exp} function, also known as the \emph{soft maximum}:
\begin{equation}
\text{LSE}_\lambda(\bu; \bq) = \lambda \log \sum_i q_i \exp(u_i / \lambda).
\label{eq:ac-adv-lse}
\end{equation}

\paragraph{Connection to Certainty Equivalents.}
The log-sum-exp is intimately related to the CES certainty equivalent. Consider the limit:
\begin{align}
\lim_{\lambda \to 0^+} \text{LSE}_\lambda(\bu; \bq) &= \max_i u_i, \label{eq:ac-adv-lse-limit0} \\
\lim_{\lambda \to \infty} \text{LSE}_\lambda(\bu; \bq) &= \sum_i q_i u_i = \mathbb{E}_q[u]. \label{eq:ac-adv-lse-limitinf}
\end{align}

The soft maximum interpolates between the hard maximum ($\lambda \to 0$) and the expectation ($\lambda \to \infty$).

Now compare with the CES certainty equivalent:
\begin{equation}
\CE_\gamma(\bc; \bq) = \left( \sum_i q_i c_i^{1-\gamma} \right)^{\frac{1}{1-\gamma}}.
\label{eq:ac-adv-ce}
\end{equation}

Taking logs:
\begin{equation}
\log \CE_\gamma = \frac{1}{1-\gamma} \log \sum_i q_i \exp\left( (1-\gamma) \log c_i \right).
\label{eq:ac-adv-ce-log}
\end{equation}

Setting $u_i = \log c_i$ and $\lambda = 1/(1-\gamma)$:
\begin{equation}
\log \CE_\gamma = \text{LSE}_{1/(1-\gamma)}(\log \bc; \bq).
\label{eq:ac-adv-ce-lse}
\end{equation}

\begin{proposition}[CES-KL Duality]
\label{prop:ac-adv-ces-kl}
The log of the CES certainty equivalent is the value function of a KL-constrained optimization problem:
\begin{equation}
\log \CE_\gamma(\bc; \bq) = \max_{\bp} \left\{ \sum_i p_i \log c_i - \frac{1}{1-\gamma} D_{KL}(\bp \| \bq) \right\}.
\label{eq:ac-adv-ces-kl}
\end{equation}
The optimal tilting is $p_i^* \propto q_i c_i^{1-\gamma}$, which is the risk-neutral measure.
\end{proposition}

\begin{calloutnote}[Interpretation]
This result provides an information-theoretic foundation for CES preferences:
\begin{itemize}[nosep]
\item The household maximizes expected log-consumption.
\item But it faces a cost for deviating from the reference distribution $\bq$.
\item Risk aversion $\gamma$ controls this cost: higher $\gamma$ means higher cost of tilting, leading to less deviation from $\bq$.
\end{itemize}

The risk-neutral measure $p_i^* \propto q_i c_i^{1-\gamma}$ emerges as the optimal belief distortion given the information cost.
\end{calloutnote}

% --------------------------------------------------------------
\subsection{Robust Control and Model Uncertainty}
\label{subsec:ac-adv-robust}
% --------------------------------------------------------------

The KL-CES connection extends naturally to robust control, where agents face model uncertainty.

\paragraph{The Robust Decision Problem.}
An agent is uncertain about the true probability distribution over states. The reference model specifies probabilities $\bq$, but the agent entertains alternatives $\bp$ within a KL ball:
\begin{equation}
\min_{\bp: D_{KL}(\bp \| \bq) \leq \kappa} \; \sum_i p_i u_i.
\label{eq:ac-adv-robust}
\end{equation}

This is the \emph{worst-case} expected utility over models within the uncertainty set.

\paragraph{Solution.}
By the same calculus:
\begin{equation}
p_i^* \propto q_i \exp(-u_i / \theta),
\label{eq:ac-adv-robust-solution}
\end{equation}
where $\theta > 0$ is the Lagrange multiplier on the KL constraint.

The worst-case model tilts toward low-utility states---the agent is pessimistic.

\paragraph{The Value Function.}
\begin{equation}
V^{\text{robust}} = -\theta \log \sum_i q_i \exp(-u_i / \theta) = -\text{LSE}_\theta(-\bu; \bq).
\label{eq:ac-adv-robust-value}
\end{equation}

This is related to the \emph{risk-sensitive} criterion in control theory and to \emph{multiplier preferences} in economics (Hansen-Sargent).

\paragraph{Connection to CES.}
With $u_i = \log c_i$:
\begin{equation}
V^{\text{robust}} = \log \left( \sum_i q_i c_i^{-1/\theta} \right)^{-\theta} = \log \CE_{1+1/\theta}(\bc; \bq).
\label{eq:ac-adv-robust-ces}
\end{equation}

Robust control with KL penalty is equivalent to CES certainty equivalent with risk aversion $\gamma = 1 + 1/\theta$.

\begin{calloutimportant}[Robustness and Risk Aversion]
The KL framework reveals that risk aversion can be interpreted as robustness to model misspecification:
\begin{itemize}[nosep]
\item A risk-neutral agent ($\gamma = 1$) uses the reference model $\bq$ directly.
\item A risk-averse agent ($\gamma > 1$) implicitly considers worst-case models, tilted toward bad outcomes.
\item Higher risk aversion corresponds to a smaller KL ball---less tolerance for model uncertainty.
\end{itemize}

This ``dual'' interpretation of risk aversion---as either curvature of utility or as robustness to misspecification---is a deep insight from the Hansen-Sargent research program.
\end{calloutimportant}

% --------------------------------------------------------------
\subsection{Extreme Value Theory}
\label{subsec:ac-adv-evt}
% --------------------------------------------------------------

We now turn to extreme value theory (EVT), which characterizes the distribution of maxima (or minima) of random variables.

\paragraph{The Setup.}
Let $X_1, X_2, \ldots, X_n$ be i.i.d. random variables with distribution $F$. Define:
\begin{equation}
M_n = \max\{X_1, \ldots, X_n\}.
\label{eq:ac-adv-Mn}
\end{equation}

The question: what is the limiting distribution of $M_n$ as $n \to \infty$?

\paragraph{The Gnedenko-Fisher-Tippett Theorem.}
If there exist normalizing sequences $a_n > 0$ and $b_n$ such that:
\begin{equation}
\frac{M_n - b_n}{a_n} \xrightarrow{d} G,
\label{eq:ac-adv-gft}
\end{equation}
then $G$ must be one of three types:
\begin{enumerate}[label=(\roman*)]
\item \textbf{Gumbel} (Type I): $G(x) = \exp(-e^{-x})$, $x \in \mathbb{R}$.
\item \textbf{Fréchet} (Type II): $G(x) = \exp(-x^{-\alpha})$, $x > 0$, $\alpha > 0$.
\item \textbf{Weibull} (Type III): $G(x) = \exp(-(-x)^{\alpha})$, $x < 0$, $\alpha > 0$.
\end{enumerate}

These can be unified into the \emph{Generalized Extreme Value} (GEV) distribution:
\begin{equation}
G_\xi(x) = \exp\left( -(1 + \xi x)^{-1/\xi} \right), \qquad 1 + \xi x > 0,
\label{eq:ac-adv-gev}
\end{equation}
where $\xi$ is the shape parameter: $\xi > 0$ (Fréchet), $\xi = 0$ (Gumbel), $\xi < 0$ (Weibull).

\paragraph{Connection to Production with Complements.}
Recall the Leontief production function:
\begin{equation}
Y = \min\{a_1 k_1, a_2 k_2, \ldots, a_n k_n\}.
\label{eq:ac-adv-leontief}
\end{equation}

If the productivities $a_i$ are random, then output $Y$ is the \emph{minimum} of $n$ random variables. By EVT, the distribution of $Y$ (suitably normalized) converges to a GEV distribution as $n \to \infty$.

\paragraph{The CES Limit.}
The CES production function:
\begin{equation}
Y = \left( \sum_{i=1}^{n} \alpha_i a_i^{1-\sigma} k_i^{\sigma} \right)^{1/\sigma}
\label{eq:ac-adv-ces-prod}
\end{equation}
interpolates between:
\begin{itemize}[nosep]
\item Sum (perfect substitutes): $\sigma = 1$, $Y = \sum_i \alpha_i a_i k_i$.
\item Geometric mean: $\sigma = 0$, $Y = \prod_i (a_i k_i)^{\alpha_i}$.
\item Minimum (perfect complements): $\sigma \to -\infty$, $Y = \min_i \{a_i k_i / \alpha_i^{-1/(1-\sigma)}\}$.
\end{itemize}

As $\sigma \to -\infty$, the CES converges to the Leontief, and EVT governs the distribution of output.

% --------------------------------------------------------------
\subsection{Pareto Distributions and Closed-Form Solutions}
\label{subsec:ac-adv-pareto}
% --------------------------------------------------------------

\paragraph{Pareto Shocks.}
Suppose each productivity $a_i$ has a Pareto distribution:
\begin{equation}
\mathbb{P}[a_i > x] = \left( \frac{\underline{a}_i}{x} \right)^{\alpha_i}, \qquad x \geq \underline{a}_i,
\label{eq:ac-adv-pareto}
\end{equation}
where $\alpha_i > 0$ is the tail exponent (higher $\alpha_i$ means thinner tail, more ``reliable'').

\paragraph{The Leontief Case.}
With Leontief production ($\sigma \to -\infty$) and full depreciation ($\delta = 1$):
\begin{equation}
Y = \min_i \{ a_i k_i \} = \min_i \{ a_i \phi_i \bar{I} \},
\label{eq:ac-adv-Y-leontief}
\end{equation}
where $\phi_i = k_i / \bar{I}$ are portfolio shares.

The distribution of $Y$ given portfolio $\boldsymbol{\phi}$ and total investment $\bar{I}$:
\begin{equation}
\mathbb{P}[Y > y] = \mathbb{P}\left[ a_i \phi_i \bar{I} > y, \; \forall i \right] = \prod_i \left( \frac{\phi_i \bar{I} \underline{a}_i}{y} \right)^{\alpha_i}.
\label{eq:ac-adv-Y-dist}
\end{equation}

This simplifies to:
\begin{equation}
\mathbb{P}[Y > y] = y^{-\sum_i \alpha_i} \cdot \prod_i \left( \phi_i \bar{I} \underline{a}_i \right)^{\alpha_i}.
\label{eq:ac-adv-Y-dist-simple}
\end{equation}

Output $Y$ has a Pareto distribution with tail exponent $\sum_i \alpha_i$.

\paragraph{The Certainty Equivalent.}
With risk aversion $\gamma$, the certainty equivalent is:
\begin{equation}
\CE = \left( \mathbb{E}[Y^{1-\gamma}] \right)^{\frac{1}{1-\gamma}}.
\label{eq:ac-adv-CE-pareto}
\end{equation}

For Pareto-distributed $Y$, this integral is tractable (assuming $\sum_i \alpha_i > 1 - \gamma$ for convergence):
\begin{equation}
\CE \propto \prod_i \left( \phi_i \underline{a}_i \right)^{\frac{\alpha_i}{\sum_j \alpha_j}}.
\label{eq:ac-adv-CE-closed}
\end{equation}

\paragraph{Optimal Portfolio.}
Maximizing $\CE$ over $\boldsymbol{\phi}$ subject to $\sum_i \phi_i = 1$:
\begin{equation}
\max_{\boldsymbol{\phi}} \prod_i \phi_i^{\alpha_i / \sum_j \alpha_j} \quad \text{s.t.} \quad \sum_i \phi_i = 1.
\label{eq:ac-adv-portfolio-pareto}
\end{equation}

This is a Cobb-Douglas optimization problem! The solution is:
\begin{equation}
\boxed{\phi_i^* = \frac{\alpha_i}{\sum_j \alpha_j}.}
\label{eq:ac-adv-phi-pareto}
\end{equation}

\begin{calloutimportant}[Allocate According to Reliability]
The optimal portfolio allocates investment proportionally to the \emph{tail exponents} $\alpha_i$:
\begin{itemize}[nosep]
\item Higher $\alpha_i$ means the right tail decays faster---capital type $i$ is more ``reliable'' (less prone to extremely low productivity).
\item Allocate more to reliable capital types.
\end{itemize}

This result is independent of risk aversion $\gamma$ (as long as the integral converges) and independent of the scale parameters $\underline{a}_i$. Only the tail behavior matters.
\end{calloutimportant}

\paragraph{Universality Conjecture.}
The Pareto case yields a remarkably simple result: $\phi_i^* = \alpha_i / \sum_j \alpha_j$. We conjecture that similar results hold more broadly:

\begin{quote}
\textit{For a wide class of distributions in the domain of attraction of the Fréchet distribution, the optimal portfolio under Leontief technology allocates according to tail exponents.}
\end{quote}

This conjecture is motivated by the GFT theorem: if the productivities are in the Fréchet domain, their minimum (which is the maximum of $-a_i$) has a limiting Fréchet distribution, and the tail exponents govern the limiting behavior.

% --------------------------------------------------------------
\subsection{Large Deviations and Rare Events}
\label{subsec:ac-adv-large-dev}
% --------------------------------------------------------------

Large deviation theory studies the probability of rare events---outcomes far from the mean.

\paragraph{The Large Deviation Principle.}
Let $S_n = \frac{1}{n} \sum_{i=1}^{n} X_i$ be the sample mean of i.i.d. random variables. By the law of large numbers, $S_n \to \mu = \mathbb{E}[X]$. Large deviation theory characterizes the rate at which $\mathbb{P}[S_n \approx x]$ decays for $x \neq \mu$:
\begin{equation}
\mathbb{P}[S_n \approx x] \asymp e^{-n I(x)},
\label{eq:ac-adv-ldp}
\end{equation}
where $I(x)$ is the \emph{rate function}.

\paragraph{The Rate Function.}
For many distributions, the rate function is the Legendre transform of the log-moment generating function:
\begin{equation}
I(x) = \sup_{\theta} \left\{ \theta x - \log \mathbb{E}[e^{\theta X}] \right\}.
\label{eq:ac-adv-rate}
\end{equation}

The rate function $I(x)$ measures how ``surprising'' the outcome $x$ is: higher $I(x)$ means $x$ is exponentially less likely.

\paragraph{Connection to KL Divergence.}
For discrete distributions, the rate function is the KL divergence:
\begin{equation}
I(\bp) = D_{KL}(\bp \| \bq),
\label{eq:ac-adv-rate-kl}
\end{equation}
where $\bp$ is the empirical distribution and $\bq$ is the true distribution.

This connects large deviations to our earlier discussion: the probability of observing an empirical distribution $\bp$ decays exponentially in $n \cdot D_{KL}(\bp \| \bq)$.

\paragraph{Implications for Portfolio Choice.}
In a portfolio context, rare events are extreme losses (or gains). The large deviation rate function governs:
\begin{itemize}[nosep]
\item The probability of portfolio values falling below a threshold.
\item The tail behavior of returns.
\item The ``value at risk'' and related measures.
\end{itemize}

With CES certainty equivalents, the connection to KL divergence implies that the agent implicitly considers the large deviation behavior of outcomes.

% --------------------------------------------------------------
\subsection{The Central Limit Theorem and Aggregation}
\label{subsec:ac-adv-clt}
% --------------------------------------------------------------

The CLT governs the distribution of sums; EVT governs maxima/minima. The CES aggregator interpolates between these extremes.

\paragraph{The CES as Interpolation.}
Consider:
\begin{equation}
M_\rho(\bx) = \left( \frac{1}{n} \sum_{i=1}^{n} x_i^\rho \right)^{1/\rho}.
\label{eq:ac-adv-M-rho}
\end{equation}

\begin{itemize}[nosep]
\item $\rho = 1$: Arithmetic mean $\bar{x} = \frac{1}{n} \sum_i x_i$. Governed by CLT.
\item $\rho \to 0$: Geometric mean $(\prod_i x_i)^{1/n}$. Log is governed by CLT.
\item $\rho \to \infty$: Maximum $\max_i x_i$. Governed by EVT.
\item $\rho \to -\infty$: Minimum $\min_i x_i$. Governed by EVT.
\end{itemize}

\paragraph{A Unified Limit Theory?}
The CES power $\rho$ determines which limit theory applies:
\begin{center}
\begin{tabular}{@{}lll@{}}
\toprule
\textbf{Range of $\rho$} & \textbf{Aggregator Behavior} & \textbf{Limit Theory} \\
\midrule
$\rho > 0$ (large) & Close to maximum & EVT (GEV) \\
$\rho \approx 1$ & Close to arithmetic mean & CLT (Gaussian) \\
$\rho \approx 0$ & Close to geometric mean & CLT for logs \\
$\rho < 0$ (large $|\rho|$) & Close to minimum & EVT (GEV) \\
\bottomrule
\end{tabular}
\end{center}

\paragraph{Conjecture: Generalized Limit Theory.}
We conjecture that there exists a ``generalized limit theory'' for CES aggregators that:
\begin{enumerate}[nosep]
\item Reduces to CLT when $\rho \approx 1$.
\item Reduces to EVT when $|\rho| \to \infty$.
\item Provides a continuous interpolation for intermediate $\rho$.
\end{enumerate}

Such a theory would characterize the distribution of CES aggregates as $n \to \infty$ and provide tools for approximating portfolio returns, production functions, and certainty equivalents with many components.

% --------------------------------------------------------------
\subsection{Entropy and the CES}
\label{subsec:ac-adv-entropy}
% --------------------------------------------------------------

\paragraph{Rényi Entropy.}
The Rényi entropy of order $\alpha$ for distribution $\bp$ is:
\begin{equation}
H_\alpha(\bp) = \frac{1}{1-\alpha} \log \sum_i p_i^\alpha.
\label{eq:ac-adv-renyi}
\end{equation}

Special cases:
\begin{itemize}[nosep]
\item $\alpha \to 1$: Shannon entropy $H_1 = -\sum_i p_i \log p_i$.
\item $\alpha = 2$: Collision entropy $H_2 = -\log \sum_i p_i^2$.
\item $\alpha \to \infty$: Min-entropy $H_\infty = -\log \max_i p_i$.
\end{itemize}

\paragraph{Connection to CES.}
The CES aggregator can be written in terms of Rényi entropy:
\begin{equation}
\log M_\rho(\bx; \bp) = \frac{1}{\rho} \log \sum_i p_i x_i^\rho = \log \left( \sum_i p_i x_i^\rho \right)^{1/\rho}.
\label{eq:ac-adv-ces-renyi}
\end{equation}

If we set $x_i = 1/p_i$ (the ``surprise'' of outcome $i$):
\begin{equation}
\log M_\rho(1/\bp; \bp) = \frac{1}{\rho} \log \sum_i p_i^{1-\rho} = \frac{1-\rho}{\rho} H_{1-\rho}(\bp).
\label{eq:ac-adv-ces-entropy}
\end{equation}

The CES of inverse probabilities is (a transform of) Rényi entropy.

\begin{calloutnote}[Information-Theoretic Interpretation]
This connection suggests that CES aggregation has an information-theoretic interpretation:
\begin{itemize}[nosep]
\item The CES curvature $\rho$ corresponds to the Rényi order $\alpha = 1 - \rho$.
\item Risk aversion in certainty equivalents corresponds to entropy measures.
\item Substitution elasticity in preferences corresponds to information diversity.
\end{itemize}

Exploring these connections may yield insights into the ``meaning'' of CES parameters beyond their standard economic interpretations.
\end{calloutnote}

% --------------------------------------------------------------
\subsection{Open Questions}
\label{subsec:ac-adv-open}
% --------------------------------------------------------------

We conclude with open questions for future research:

\begin{enumerate}[label=\textbf{Q\arabic*.}]
\item \textbf{Generalized limit theory}: Can we develop a unified limit theory for CES aggregates that interpolates between CLT and EVT?

\item \textbf{Optimal portfolios under EVT}: For general distributions in the domain of attraction of GEV, what is the optimal portfolio under Leontief (or near-Leontief) production? Does the ``allocate by tail exponent'' rule extend?

\item \textbf{Information-theoretic foundations}: Can CES preferences be axiomatized from information-theoretic principles (e.g., maximizing utility subject to information costs)?

\item \textbf{Large deviations and tail risk}: How does the CES certainty equivalent relate to tail risk measures like Value-at-Risk or Expected Shortfall? Can the KL connection provide new insights?

\item \textbf{Aggregation with dependence}: The results above assume independence. How do correlations between shocks affect the CES-EVT-KL connections?

\item \textbf{Dynamic extensions}: How do these information-theoretic and limit-theory connections extend to dynamic settings with recursive preferences?
\end{enumerate}

% --------------------------------------------------------------
\subsection{Summary}
\label{subsec:ac-adv-summary}
% --------------------------------------------------------------

\begin{table}[H]
\centering
\caption{Connections to Information Theory and Extreme Values}
\label{tab:ac-adv-summary}
\renewcommand{\arraystretch}{1.5}
\begin{tabular}{@{}lll@{}}
\toprule
\textbf{Area} & \textbf{Key Object} & \textbf{CES Connection} \\
\midrule
Information theory & KL divergence $D_{KL}(\bp \| \bq)$ & Risk aversion = info cost \\
Robust control & Worst-case model & CES certainty equivalent \\
Extreme value theory & GEV distribution & Leontief $\leftrightarrow$ min \\
Pareto distributions & Tail exponent $\alpha_i$ & Optimal portfolio $\phi_i^* = \alpha_i / \sum \alpha_j$ \\
Large deviations & Rate function $I(x)$ & KL divergence \\
Entropy & Rényi entropy $H_\alpha$ & CES curvature $\rho = 1 - \alpha$ \\
\bottomrule
\end{tabular}
\end{table}

\begin{callouttip}[Key Insights]
The speculative connections in this section suggest:
\begin{itemize}[nosep]
\item \textbf{CES is fundamental}: The CES structure emerges from information-theoretic principles (KL divergence) and from limit theorems (EVT, CLT).
\item \textbf{Risk aversion = information cost}: High risk aversion corresponds to high cost of deviating from reference beliefs.
\item \textbf{Tail exponents matter}: For Leontief production with Pareto shocks, optimal portfolios allocate by reliability (tail exponents).
\item \textbf{Unified framework}: CES interpolates between CLT (sums) and EVT (extremes), suggesting a generalized limit theory.
\item \textbf{Entropy connection}: CES curvature maps to Rényi entropy order, hinting at information-theoretic foundations.
\end{itemize}
These connections merit further investigation and may reveal deeper structure underlying the CES framework.
\end{callouttip}
